\chapter{Contexto y estado del arte}
\label{cap:contexto}

La máquina, los materiales, el flujo de pedidos y las condiciones reales de captura en Innokey determinan qué técnicas resultan útiles en el taller. El capítulo revisa el registro geométrico, los métodos de segmentación y la conversión de máscaras en trayectorias de corte, y sitúa el papel de Random Forest y de los modelos fundacionales dentro de la comparación experimental.

\section{Contexto productivo}

La pieza, el flujo CAD/CAM y las condiciones de captura se describen primero porque de ellos nacen las restricciones del procesamiento.

\subsection{Fabricación personalizada de toppers}

La fabricación por encargo obliga a preparar la máquina muchas veces para lotes pequeños. Cada pedido puede incluir formas distintas, tamaños variables y materiales parcialmente aprovechados. En Innokey, los toppers muestran con claridad esa variación. Los diseños parten de imágenes o ideas entregadas por los clientes y luego se corrigen, depuran o redibujan antes de imprimirlos sobre una lámina adhesiva.

El punto crítico aparece al integrar la impresión y la base rígida. El plóter recorta la capa adhesiva y el láser corta el MDF o el acrílico en operaciones separadas. Más tarde, el operario debe hacer coincidir cada par de piezas. La diversidad de formas vuelve poco prácticas las guías rígidas y una desviación pequeña puede dejar visible el soporte.

\subsection{Flujo CAD/CAM y sistema de coordenadas de trabajo}

Las cortadoras láser CNC trabajan con trayectorias vectoriales planas. La geometría suele prepararse mediante diseño asistido por computador (\textit{computer-aided design}, CAD) y después pasa al software de fabricación asistida por computador (\textit{computer-aided manufacturing}, CAM), donde se establecen velocidad, potencia y orden de corte. Existen antecedentes que convierten fotografías o imágenes rasterizadas en dimensiones, geometría vectorial y trayectorias compatibles con fabricación \cite{Bhandari2023,Jednorg2025,Khafagy2025}. El registro de una impresión ya adherida a su soporte constituye la aplicación específica abordada en este trabajo.

Conviene distinguir el sistema de coordenadas de máquina del sistema de coordenadas de trabajo. El primero pertenece a la estructura de la cortadora y a sus finales de carrera. El segundo, denominado WCS por \textit{Work Coordinate System}, lo define el operario según el material y el aprovechamiento de la mesa. En una máquina de aproximadamente $130 \times 90\,\text{cm}$, ese origen puede cambiar entre pedidos. Por este motivo, la herramienta no depende de una plantilla fija. Utiliza una marca en forma de L grabada por el propio láser sobre una zona libre de la hoja.

\subsection{Captura con cámara de consumo}

La región de trabajo se limita a una hoja A4 adherida al soporte. La restricción concentra los píxeles de la cámara en el área que realmente contiene los diseños y permite trabajar a una escala nominal de $10\,\text{px/mm}$. Cubrir toda la mesa exigiría alejar la cámara y reduciría el detalle disponible para los bordes finos.

La captura se realiza con un teléfono inteligente montado en la compuerta superior. El soporte es removible y no garantiza una posición idéntica entre sesiones. La hoja incorpora cuatro marcadores negros que permiten estimar su orientación y rectificar la perspectiva. Esta decisión reduce el coste del montaje y hace posible reproducirlo sin una cámara industrial fija.

\subsection{Dificultades visuales de la escena}

La impresión sobre fondo blanco facilita parte de la separación, pero no elimina la dificultad. Los diseños contienen tonos pastel, zonas blancas internas, degradados y trazos oscuros que no pertenecen al borde de corte. La iluminación del taller añade sombras, reflejos y cambios de balance de blancos. También pueden aparecer en la captura el cabezal y la guía del eje X.

El sistema debe recuperar la silueta externa sin incorporar al contorno de corte rasgos gráficos internos, como ojos, pliegues o detalles decorativos. La Figura~\ref{fig:dificultades_visuales} reúne varias de las condiciones observadas durante la adquisición.

\figuraopcionaldetallada{dificultades_visuales_toppers}{0.9\textwidth}{Dificultades visuales presentes en las capturas reales.}{Dificultades visuales presentes en las capturas reales. El panel superior izquierdo muestra una iluminación desigual, con una sombra horizontal que atraviesa la hoja y el soporte. El panel superior derecho incorpora parte del cabezal dentro del campo visual. En la fila inferior se observan, respectivamente, la rotación de la hoja y su desplazamiento lateral respecto a la posición habitual. Los marcadores de las esquinas permiten rectificar la escena antes de segmentar los diseños pese a esas diferencias.}{fig:dificultades_visuales}

\section{Fundamentos geométricos}

El registro de la impresión exige distinguir dos problemas geométricos. La calibración describe cómo la cámara forma la imagen, mientras que la homografía relaciona la hoja observada con un plano de dimensiones conocidas.

\subsection{Calibración y distorsión de cámara}

La calibración intrínseca estima la distancia focal, el punto principal y los coeficientes que describen la distorsión del objetivo. El procedimiento de Zhang sigue siendo una referencia para obtener estos parámetros a partir de varias observaciones de un patrón plano \cite{Zhang2000}. Trabajos posteriores muestran que la calidad de las esquinas detectadas y las diferencias entre zonas del sensor influyen en el resultado de la calibración \cite{Wang2022}.

El montaje se evaluó sin completar una calibración intrínseca formal: la hoja ocupa la zona central del encuadre y sus márgenes no mostraron una curvatura visible durante la inspección, de modo que la perspectiva se corrigió mediante una transformación planar. Esa elección resuelve el registro de la hoja, pero no cuantifica la distorsión del objetivo ni demuestra que sea despreciable. La posible deformación radial residual se conserva, por tanto, como una amenaza a la validez geométrica.

\subsection{Homografía y rectificación planar}

Una homografía es una transformación proyectiva entre dos planos. Permite relacionar la hoja observada en perspectiva con un rectángulo de dimensiones conocidas. En coordenadas homogéneas, un punto $(x,y)$ de la imagen original se transforma en $(x',y')$ mediante

\begin{equation}
\begin{bmatrix}
x' \\ y' \\ 1
\end{bmatrix}
\sim
H
\begin{bmatrix}
x \\ y \\ 1
\end{bmatrix},
\qquad
H =
\begin{bmatrix}
h_{11} & h_{12} & h_{13} \\
h_{21} & h_{22} & h_{23} \\
h_{31} & h_{32} & 1
\end{bmatrix}.
\label{eq:homografia}
\end{equation}

Las cuatro correspondencias de esquina aportadas por los marcadores permiten estimar $H$ y llevar la hoja a un lienzo de $2100 \times 2970$ píxeles. Investigaciones recientes confirman la utilidad de la homografía para corregir distorsión de perspectiva en mediciones planas \cite{Wang2025}. La Figura~\ref{fig:homografia_contexto} muestra el sentido de esta transformación dentro del sistema.

\figuraopcional{fig_homografia}{0.88\textwidth}{Rectificación de la hoja mediante correspondencias entre los marcadores detectados y el plano canónico A4.}{fig:homografia_contexto}

Los marcadores fiduciales se utilizan con frecuencia para recuperar posición y orientación a partir de patrones conocidos \cite{GarridoJurado2014}. Aquí los cuatro cuadrados delimitan la hoja, mientras que la marca en L cumple una función diferente. Su vértice define el origen WCS y sus brazos aportan la dirección de los ejes.

\section{Segmentación mediante visión clásica}

Los métodos clásicos separan regiones a partir de propiedades expresadas de forma explícita. En este caso resultan útiles la saturación, la luminosidad, los gradientes y la conectividad espacial. La corrección de zonas sobreexpuestas y la segmentación adaptativa se han empleado en escenas industriales con reflejos \cite{Li2023}, mientras que los operadores adaptativos de borde buscan conservar límites de pieza bajo ruido variable \cite{Li2024}.

El detector de Canny combina suavizado, gradiente, supresión de no máximos e histéresis \cite{Canny1986}. En una imagen de \textit{topper} también responde a numerosas líneas internas, por lo que no se emplea de forma aislada. Los bordes detectados se combinan con umbrales en HSV y operaciones morfológicas que cierran discontinuidades y rellenan la región externa. Así se obtiene una línea base interpretable que, además, localiza de forma aproximada las cajas que recibe SAM 2.

\section{Aprendizaje automático para segmentación}

Los métodos de aprendizaje considerados representan dos niveles de complejidad. Random Forest aprende una decisión por píxel a partir de características diseñadas, mientras que SAM 2 reutiliza una representación neuronal preentrenada y recibe una guía espacial mediante prompts.

Los antecedentes industriales citados permiten concretar qué significa esa incorporación de IA. Ozkan compara regresión logística, máquinas de vectores de soporte, árboles de decisión, Random Forest, XGBoost y una red neuronal artificial después de extraer descriptores geométricos. La red neuronal obtuvo el mejor resultado en la detección de defectos de casquillos metálicos \cite{Ozkan2025}. Loganathan et al. emplean YOLOv8n-seg para segmentación multiclase de defectos de soldadura y medición geométrica en tiempo real \cite{Loganathan2026}. Ambos trabajos combinan aprendizaje con procesamiento visual previo y muestran que el modelo debe describirse por su función concreta dentro de la cadena.

\subsection{Random Forest como clasificador por píxel}

Random Forest agrega las predicciones de varios árboles entrenados con muestreo y selección aleatoria de características \cite{Breiman2001}. La combinación reduce la varianza de un árbol individual y permite modelar relaciones no lineales sin una red neuronal profunda. En una segmentación por píxel, cada observación se describe mediante color, gradiente y textura local. El modelo estima entonces la probabilidad de que ese píxel pertenezca al topper.

\figuraopcional{fig_random_forest}{0.9\textwidth}{Funcionamiento del Random Forest por píxel, desde la extracción de características hasta la agregación de los árboles y la formación de la máscara.}{fig:rf_contexto}

La Figura~\ref{fig:rf_contexto} resume el procedimiento. Para este proyecto, Random Forest ofrece una referencia supervisada distinta de las reglas clásicas y puede entrenarse con un volumen moderado de datos. La contrapartida es su dependencia del dominio. El clasificador puede aprender con gran precisión el blanco digital, las sombras programadas y las texturas sintéticas, pero fallar cuando la cámara introduce ruido, compresión y variaciones cromáticas ausentes durante el entrenamiento.

\subsection{Redes convolucionales y alternativas consideradas}

Las redes convolucionales aprenden representaciones jerárquicas a partir de los datos \cite{LeCun1998}. U-Net combina un codificador con un decodificador y conexiones de salto para recuperar detalle espacial \cite{Ronneberger2015}. Mask R-CNN añade una rama de máscara a un detector regional \cite{He2017}, mientras que YOLO plantea una detección directa orientada a velocidad \cite{Redmon2016}.

U-Net, Mask R-CNN y YOLO podrían aplicarse al problema, pero entrenarlas desde cero o compararlas mediante un ajuste fino sólido exigiría más diversidad real anotada. Las trece fotografías disponibles corresponden a una sola lámina física, de modo que una red propia produciría una evaluación poco defendible. El banco sintético ayuda a estudiar la transferencia, aunque no sustituye la variabilidad de producción. Por ello se optó por un modelo preentrenado capaz de trabajar con prompts, sin necesidad de ajustarlo con las imágenes del proyecto.

\subsection{Modelos fundacionales y SAM 2}

Para este problema, el interés de un modelo fundacional reside en que no parte de las trece fotografías del taller para aprender qué aspecto tiene un objeto. Aprovecha representaciones adquiridas durante un preentrenamiento amplio y las orienta a la tarea mediante prompts, ajuste fino o módulos ligeros. En segmentación, SAM mostró que puntos o cajas podían guiar la producción de máscaras incluso sin un entrenamiento específico para cada distribución visual \cite{Kirillov2023}.

SAM 2 extiende la segmentación condicionada a imágenes y vídeo mediante una arquitectura que combina un codificador visual jerárquico, un codificador de prompts y un decodificador de máscaras. La memoria temporal es relevante para vídeo, aunque en este TFM cada hoja se procesa como una imagen independiente \cite{Ravi2025}. De las variantes disponibles, este prototipo emplea Hiera Tiny, la de menor tamaño, ejecutada en CPU, cuyo coste se cuantifica mediante mediciones propias.

\figuraopcional{fig_sam2}{0.94\textwidth}{Arquitectura funcional de SAM 2 para una imagen estática condicionada mediante cajas geométricas.}{fig:sam2_contexto}

La Figura~\ref{fig:sam2_contexto} diferencia la imagen, el prompt y la máscara resultante. En este trabajo se carga un checkpoint preentrenado y se comparan distintas formas de construir el prompt, sin ajustar SAM 2 con los datos del TFM. Así, lo que se evalúa es la transferencia de un modelo fundacional, no un entrenamiento realizado dentro del proyecto. Esta configuración tiene una contrapartida: SAM 2 necesita una localización inicial. En las capturas reales, las cajas proceden de las componentes obtenidas por el módulo clásico, aunque sus bordes no se reutilizan como máscara final.

\section{Métricas de segmentación}

La evaluación de una máscara exige medir algo más que el porcentaje de píxeles acertados, porque el fondo ocupa gran parte de la hoja. La intersección sobre unión, conocida como IoU o índice de Jaccard, divide el área común entre la unión de la predicción y la referencia. El coeficiente Dice duplica la intersección y la divide entre la suma de ambas áreas. Las dos medidas toman el valor uno cuando las máscaras coinciden y disminuyen con los falsos positivos y falsos negativos. Dice suele producir un valor numérico mayor que IoU para el mismo par de máscaras, aunque ambas ordenan de forma coherente el grado de solapamiento \cite{Taha2015}.

El Boundary F1 traslada la comparación al borde. Un punto predicho se considera correcto cuando cae dentro de una tolerancia alrededor del contorno de referencia. A partir de la precisión y la exhaustividad se calcula su media armónica. La implementación empleada es una medida F1 de contorno con tolerancia de tres píxeles, inspirada en la evaluación tolerante de límites de Martin et al. \cite{Martin2004}. No reproduce el emparejamiento uno a uno del protocolo BSDS. Esta métrica resulta útil para corte porque dos máscaras pueden tener un área similar y, sin embargo, diferir en apéndices finos o entrantes de la silueta. El error de componentes complementa el análisis al comparar la cantidad de regiones conectadas con las instancias esperadas.

\section{Vectorización y geometría fabricable}

Una máscara binaria todavía no constituye una trayectoria de máquina. El contorno contiene una gran cantidad de vértices y pequeñas irregularidades propias del raster. Una polilínea es una secuencia ordenada de segmentos rectos que aproxima ese límite y que puede cerrarse para describir la silueta completa.

El algoritmo de Douglas-Peucker reduce los puntos redundantes dentro de una tolerancia geométrica $\varepsilon$ \cite{DouglasPeucker1973}. El procedimiento une primero los extremos de una secuencia y localiza el punto con mayor distancia perpendicular al segmento. Si esa distancia supera $\varepsilon$, divide la curva en ese punto y repite el cálculo en cada parte. Si no lo supera, conserva únicamente los extremos. Un valor pequeño mantiene más detalle y produce más vértices, mientras que uno grande simplifica la trayectoria y puede deformar rasgos finos. Por ello, la tolerancia debe expresarse en relación con la escala física de la imagen.

La literatura que conecta imagen y fabricación confirma la necesidad de validar la salida como geometría y no solo como visualización. Bhandari y Manandhar combinan visión y CAD para recuperar dimensiones \cite{Bhandari2023}. Jednoróg desarrolla una herramienta orientada a generar DXF desde fotografías en un entorno de pequeña empresa \cite{Jednorg2025}. Khafagy et al. estudian la extracción y optimización de trayectorias complejas desde gráficos vectoriales para aplicaciones robóticas \cite{Khafagy2025}. En todos los casos, la utilidad final depende de conservar escala, conectividad y referencia espacial.

\section{Síntesis crítica}

La revisión respalda una arquitectura híbrida con responsabilidades diferenciadas. La geometría controlada se ocupa de la rectificación y del WCS, lo que permite auditar el registro. El módulo clásico propone cajas aproximadas y SAM 2 trabaja sobre esas regiones para generar la máscara que se vectoriza. Random Forest aporta una referencia supervisada con la que estudiar el aprendizaje por píxel y la brecha entre los dominios sintético y real.

Cada componente se juzga según la tarea que resuelve, no por la expectativa de que la inteligencia artificial mejore todas las métricas. Una regla ajustada al entorno puede conservar mayor precisión, mientras un modelo preentrenado puede sostener la segmentación cuando cambia la apariencia. La aplicación usa SAM 2 de forma efectiva y trazable, y mantiene visibles los componentes geométricos que conectan la imagen con la máquina.
